\section{چکیده}

در این آزمایش، رفتار سه پروتکل مسیریابی پویای \لر{OSPF}، \لر{EIGRP} و
\لر{RIPv2} روی یک توپولوژی یکسان در \لر{Packet Tracer} بررسی شد.
پیکربندی همسایگی‌ها و جدول‌های مسیریابی، ساختار بسته‌های کنترلی، تاخیر
\لر{ping}، مسیر پیش‌فرض، جداسازی \لر{Internet} از اطلاعات داخلی، توازن بار،
اعمال ترجیح مسیر در \لر{R4} و واکنش شبکه به خاموش شدن \لر{R2} ارزیابی شد.
نتیجه‌ها نشان دادند معیار \لر{OSPF} بر هزینه پیوند، معیار \لر{EIGRP} بر متریک
ترکیبی پهنای باند و تاخیر و معیار \لر{RIP} بر تعداد \لر{hop} استوار است.
همچنین محدودیت \لر{Packet Tracer 8.2.1} در نمایش توازن بار دقیق به‌ازای هر
بسته مشخص شد و روش انجام آن روی تجهیزات واقعی تشریح گردید.

\بدون‌تورفتگی \مهم{واژه‌های کلیدی}:
\لر{Packet Tracer}، \لر{OSPF}، \لر{EIGRP}، \لر{RIPv2}، مسیریابی پویا،
همگرایی، مسیر پیش‌فرض، \لر{NAT}، توازن بار

\section{۱. پیاده‌سازی توپولوژی و آدرس‌دهی}

توپولوژی آزمایش چهارم مبنای کار قرار گرفت و با شماره گروه ۶ تنظیم شد.
شبکه انتهایی \لر{PC1} برابر \لر{10.10.6.0/24} و شبکه انتهایی \لر{PC3}
برابر \لر{10.10.16.0/24} است. پیوندهای میانی به‌ترتیب از شبکه‌های
\لر{10.10.8.0/24} برای \لر{R1-R2}، \لر{10.10.9.0/24} برای \لر{R2-R3}،
\لر{10.10.10.0/24} برای \لر{R3-R4}، \لر{10.10.11.0/24} برای
\لر{R1-R5} و \لر{10.10.12.0/24} برای \لر{R4-R5} استفاده می‌کنند.
نشانی \لر{PC1} برابر \لر{10.10.6.1} با دروازه \لر{10.10.6.2} و نشانی
\لر{PC3} برابر \لر{10.10.16.1} با دروازه \لر{10.10.16.2} قرار گرفت.
واسط بیرونی \لر{R5} برابر \لر{213.80.11.4/24} و واسط روتر
\لر{Internet} برابر \لر{213.80.11.5/24} تنظیم شد.

برای یکسان ماندن شرایط مقایسه، فایل پایه در سه نسخه مستقل
\لر{OSPF.pkt}، \لر{EIGRP.pkt} و \لر{RIP.pkt} ذخیره شد. در هر نسخه فقط
پروتکل مربوط به همان فایل روی \لر{R1} تا \لر{R5} فعال است و روتر
\لر{Internet} در هیچ فرایند مسیریابی پویا شرکت ندارد.

\section{۲. راه‌اندازی \لر{OSPF}، \لر{EIGRP} و \لر{RIP}}

\subsection{پیکربندی \لر{OSPF}}

در فایل \لر{OSPF}، فرایند شماره ۱ روی پنج روتر اجرا شد. همه شبکه‌های داخلی
در \لر{area 0} قرار گرفتند، برای هر روتر \لر{Router ID} مستقل تعریف شد و با
\کد{passive-interface default} ارسال بسته کنترلی روی واسط‌های انتهایی متوقف
گردید. فقط واسط‌های متصل به روترهای مجاور از حالت \لر{passive} خارج شدند.
\لر{R5} نیز مسیر پیش‌فرض را به داخل \لر{OSPF} منتشر کرد.

\شروع{شکل}[H]
\centerimg{q02-ospf-r1-configuration.png}{11cm}
\شرح{پیکربندی \لر{OSPF} روی \لر{R1} با \لر{Router ID} برابر \لر{1.1.1.1} و دو واسط فعال داخلی}
\برچسب{شکل:q02-ospf-r1-config}
\پایان{شکل}

\شروع{شکل}[H]
\centerimg{q02-ospf-r2-configuration.png}{11cm}
\شرح{پیکربندی \لر{OSPF} روی \لر{R2} و فعال بودن دو پیوند \لر{R1} و \لر{R3}}
\برچسب{شکل:q02-ospf-r2-config}
\پایان{شکل}

\شروع{شکل}[H]
\centerimg{q02-ospf-r3-configuration.png}{11cm}
\شرح{پیکربندی \لر{OSPF} روی \لر{R3} و فعال بودن پیوندهای \لر{R2} و \لر{R4}}
\برچسب{شکل:q02-ospf-r3-config}
\پایان{شکل}

\شروع{شکل}[H]
\centerimg{q02-ospf-r4-configuration.png}{11cm}
\شرح{پیکربندی \لر{OSPF} روی \لر{R4} و \لر{passive} ماندن شبکه \لر{PC3}}
\برچسب{شکل:q02-ospf-r4-config}
\پایان{شکل}

\شروع{شکل}[H]
\centerimg{q02-ospf-r5-configuration.png}{11cm}
\شرح{پیکربندی \لر{OSPF} روی \لر{R5} و انتشار مسیر پیش‌فرض}
\برچسب{شکل:q02-ospf-r5-config}
\پایان{شکل}

جدول همسایه \لر{R1} دو مجاورت \لر{FULL} با \لر{R2} و \لر{R5} را نشان
می‌دهد. در \لر{R4} نیز \لر{R3} و \لر{R5} در حالت \لر{FULL} قرار دارند.
این وضعیت تشکیل درست همسایگی‌ها و تبادل پایگاه وضعیت پیوند
\پانویس{\lr{Link-State Database}} را تایید می‌کند.

\شروع{شکل}[H]
\centerimg{q02-ospf-r1-neighbors-and-routes.png}{13.5cm}
\شرح{همسایه‌های \لر{FULL} و مسیرهای \لر{OSPF} در \لر{R1}}
\برچسب{شکل:q02-ospf-r1-neighbors}
\پایان{شکل}

\شروع{شکل}[H]
\centerimg{q02-ospf-r4-neighbors-and-routes.png}{13.5cm}
\شرح{همسایه‌های \لر{FULL} و مسیرهای \لر{OSPF} در \لر{R4}}
\برچسب{شکل:q02-ospf-r4-neighbors}
\پایان{شکل}

\subsection{پیکربندی \لر{EIGRP}}

در فایل \لر{EIGRP}، سیستم خودمختار ۱۰۰ روی \لر{R1} تا \لر{R5} اجرا شد.
خلاصه‌سازی خودکار غیرفعال و شبکه‌های \لر{10.10.0.0/16} وارد فرایند شدند.
روش \کد{passive-interface default} همان مرز دقیق واسط‌های داخلی و انتهایی را
ایجاد کرد. \لر{R5} مسیر ایستای پیش‌فرض را با \لر{seed metric} مشخص در
\لر{EIGRP} بازتوزیع کرد.

\شروع{شکل}[H]
\centerimg{q02-eigrp-r1-configuration.png}{13.5cm}
\شرح{پیکربندی \لر{EIGRP} و تشکیل دو همسایگی در \لر{R1}}
\برچسب{شکل:q02-eigrp-r1-config}
\پایان{شکل}

\شروع{شکل}[H]
\centerimg{q02-eigrp-r2-configuration.png}{13.5cm}
\شرح{پیکربندی \لر{EIGRP} و تشکیل همسایگی‌های \لر{R1} و \لر{R3} در \لر{R2}}
\برچسب{شکل:q02-eigrp-r2-config}
\پایان{شکل}

\شروع{شکل}[H]
\centerimg{q02-eigrp-r3-configuration.png}{13.5cm}
\شرح{پیکربندی \لر{EIGRP} و تشکیل همسایگی‌های \لر{R2} و \لر{R4} در \لر{R3}}
\برچسب{شکل:q02-eigrp-r3-config}
\پایان{شکل}

\شروع{شکل}[H]
\centerimg{q02-eigrp-r4-configuration.png}{13.5cm}
\شرح{پیکربندی \لر{EIGRP} و تشکیل همسایگی‌های \لر{R3} و \لر{R5} در \لر{R4}}
\برچسب{شکل:q02-eigrp-r4-config}
\پایان{شکل}

\شروع{شکل}[H]
\centerimg{q02-eigrp-r5-configuration.png}{13.5cm}
\شرح{پیکربندی \لر{EIGRP}، همسایگی‌های داخلی و بازتوزیع مسیر پیش‌فرض در \لر{R5}}
\برچسب{شکل:q02-eigrp-r5-config}
\پایان{شکل}

خروجی همسایه‌ها در \لر{R1} و \لر{R4} مقدار \لر{Q Cnt} برابر صفر دارد؛
بنابراین بسته‌ای در صف نمانده و مجاورت‌ها پایدار بوده‌اند. مسیرهای دارای کد
\لر{D} نیز یادگیری شبکه‌های دوردست با \لر{EIGRP} را نشان می‌دهند.

\شروع{شکل}[H]
\centerimg{q02-eigrp-r1-neighbors-and-routes.png}{13cm}
\شرح{همسایه‌ها و مسیر یادگرفته‌شده \لر{EIGRP} در \لر{R1}}
\برچسب{شکل:q02-eigrp-r1-neighbors}
\پایان{شکل}

\شروع{شکل}[H]
\centerimg{q02-eigrp-r4-neighbors-and-routes.png}{13cm}
\شرح{همسایه‌ها و مسیر یادگرفته‌شده \لر{EIGRP} در \لر{R4}}
\برچسب{شکل:q02-eigrp-r4-neighbors}
\پایان{شکل}

\subsection{پیکربندی \لر{RIP}}

در فایل \لر{RIP}، نسخه ۲ انتخاب شد تا مسیریابی بدون کلاس و ارسال
\لر{multicast} در دسترس باشد. \کد{no auto-summary} از خلاصه‌سازی در مرز
شبکه اصلی جلوگیری کرد. شبکه \لر{10.0.0.0} وارد \لر{RIP} شد و فقط پیوندهای
روتر به روتر از حالت \لر{passive} خارج شدند. \لر{R5} مسیر پیش‌فرض را با
\کد{default-information originate} منتشر کرد.

\شروع{شکل}[H]
\centerimg{q02-rip-r1-configuration.png}{11cm}
\شرح{پیکربندی \لر{RIPv2} روی \لر{R1}}
\برچسب{شکل:q02-rip-r1-config}
\پایان{شکل}

\شروع{شکل}[H]
\centerimg{q02-rip-r2-configuration.png}{11cm}
\شرح{پیکربندی \لر{RIPv2} روی \لر{R2}}
\برچسب{شکل:q02-rip-r2-config}
\پایان{شکل}

\شروع{شکل}[H]
\centerimg{q02-rip-r3-configuration.png}{11cm}
\شرح{پیکربندی \لر{RIPv2} روی \لر{R3}}
\برچسب{شکل:q02-rip-r3-config}
\پایان{شکل}

\شروع{شکل}[H]
\centerimg{q02-rip-r4-configuration.png}{11cm}
\شرح{پیکربندی \لر{RIPv2} روی \لر{R4}}
\برچسب{شکل:q02-rip-r4-config}
\پایان{شکل}

\شروع{شکل}[H]
\centerimg{q02-rip-r5-configuration.png}{11cm}
\شرح{پیکربندی \لر{RIPv2} و انتشار مسیر پیش‌فرض روی \لر{R5}}
\برچسب{شکل:q02-rip-r5-config}
\پایان{شکل}

پایگاه \لر{RIP} در \لر{R1} مسیرهای دریافت‌شده از \لر{R2} و \لر{R5} را
همراه متریک تعداد \لر{hop} نگه می‌دارد. پایگاه \لر{R4} نیز دو جهت \لر{R3}
و \لر{R5} را نشان می‌دهد. خروجی \لر{R2} دوره به‌روزرسانی ۳۰ ثانیه،
\لر{invalid timer} برابر ۱۸۰ ثانیه، \لر{hold-down} برابر ۱۸۰ ثانیه و
\لر{flush} برابر ۲۴۰ ثانیه را ثبت کرده است.

\شروع{شکل}[H]
\centerimg{q02-rip-r1-database-and-default.png}{14cm}
\شرح{پایگاه \لر{RIP} و مسیر پیش‌فرض در \لر{R1}}
\برچسب{شکل:q02-rip-r1-database}
\پایان{شکل}

\شروع{شکل}[H]
\centerimg{q02-rip-r4-database-and-default.png}{14cm}
\شرح{پایگاه \لر{RIP} و مسیرهای دریافتی از دو جهت در \لر{R4}}
\برچسب{شکل:q02-rip-r4-database}
\پایان{شکل}

\شروع{شکل}[H]
\centerimg{q02-rip-r2-timers-and-routes.png}{13cm}
\شرح{زمان‌سنج‌ها، منابع اطلاعات و مسیرهای \لر{RIP} در \لر{R2}}
\برچسب{شکل:q02-rip-r2-timers}
\پایان{شکل}

\section{۳. بررسی بسته‌های کنترلی در \لر{Simulation}}

در \لر{OSPF}، بسته \لر{Hello} با مقصد \لر{multicast} برابر
\لر{224.0.0.5} ارسال می‌شود. نمونه نخست از \لر{10.10.9.2} در \لر{R3} به
\لر{R2} رسیده و نمونه دوم از \لر{10.10.12.2} در \لر{R5} به \لر{R4}
رسیده است. سربرگ لایه سوم نوع \لر{OSPF HELLO} و سربرگ \لر{Ethernet II}
نگاشت \لر{multicast} را نشان می‌دهد. \لر{Hello} برای کشف همسایه، نگهداری
مجاورت و تطبیق پارامترهایی مانند \لر{area} و زمان‌سنج‌ها استفاده می‌شود.
پس از تشکیل مجاورت، \لر{Database Description}، \لر{Link State Request}،
\لر{Link State Update} و \لر{Link State Acknowledgment} برای همگام‌سازی
\لر{LSDB} مبادله می‌شوند.

\شروع{شکل}[H]
\centerimg{q03-ospf-hello-r3-to-r2.png}{13cm}
\شرح{بسته \لر{OSPF Hello} از \لر{R3} به \لر{R2} با مقصد \لر{224.0.0.5}}
\برچسب{شکل:q03-ospf-r3-r2}
\پایان{شکل}

\شروع{شکل}[H]
\centerimg{q03-ospf-hello-r5-to-r4.png}{13cm}
\شرح{بسته \لر{OSPF Hello} از \لر{R5} به \لر{R4} روی پیوند داخلی}
\برچسب{شکل:q03-ospf-r5-r4}
\پایان{شکل}

بسته‌های کنترلی \لر{EIGRP} از \لر{multicast} برابر \لر{224.0.0.10}
استفاده می‌کنند. دو نمونه ثبت‌شده، ترافیک \لر{EIGRP} نسخه ۲ را روی پیوند
\لر{R3-R2} و \لر{R5-R1} نشان می‌دهند. \لر{EIGRP} با \لر{Hello} همسایه
را نگه می‌دارد و \لر{Update} را فقط هنگام نیاز و به‌شکل محدود ارسال می‌کند.
الگوریتم \لر{DUAL} با نگهداری \لر{Successor} و در صورت وجود
\لر{Feasible Successor}\پانویس{\lr{Feasible Successor}}، مسیر بدون حلقه را
انتخاب می‌کند؛ \لر{Query} و \لر{Reply} نیز هنگام نبود مسیر جانشین مناسب برای
محاسبه دوباره به کار می‌روند.

\شروع{شکل}[H]
\centerimg{q03-eigrp-control-packet-r3-to-r2.png}{13cm}
\شرح{بسته کنترلی \لر{EIGRP} از \لر{R3} به \لر{R2} با مقصد \لر{224.0.0.10}}
\برچسب{شکل:q03-eigrp-r3-r2}
\پایان{شکل}

\شروع{شکل}[H]
\centerimg{q03-eigrp-control-packet-r5-to-r1.png}{13cm}
\شرح{بسته کنترلی \لر{EIGRP} از \لر{R5} به \لر{R1}}
\برچسب{شکل:q03-eigrp-r5-r1}
\پایان{شکل}

\لر{RIPv2} پیام \لر{Response} را با \لر{UDP} و پورت مبدا و مقصد ۵۲۰ به
\لر{multicast} برابر \لر{224.0.0.9} می‌فرستد. \لر{Command} برابر ۲ در هر
دو نمونه نشان‌دهنده \لر{Response} است. این پیام فهرست شبکه‌ها و متریک آن‌ها
را حمل می‌کند و به‌طور معمول هر ۳۰ ثانیه تکرار می‌شود.

\شروع{شکل}[H]
\centerimg{q03-rip-update-r3-to-r2.png}{13cm}
\شرح{پیام \لر{RIPv2 Response} از \لر{R3} به \لر{R2} با \لر{UDP/520}}
\برچسب{شکل:q03-rip-r3-r2}
\پایان{شکل}

\شروع{شکل}[H]
\centerimg{q03-rip-update-r5-to-r1.png}{13cm}
\شرح{پیام \لر{RIPv2 Response} از \لر{R5} به \لر{R1}}
\برچسب{شکل:q03-rip-r5-r1}
\پایان{شکل}

\section{۴. اندازه‌گیری و مقایسه تاخیر \لر{ping}}

برای قابل مشاهده شدن اثر متریک، پهنای باند اسمی پیوندهای میان روترها کاهش و
\لر{delay} آن‌ها افزایش یافت. با این حال \لر{Packet Tracer} پردازش شبکه محلی
را غالباً با زمان کمتر از یک میلی‌ثانیه نمایش می‌دهد؛ بنابراین نتیجه‌های
ثبت‌شده بر مقدار واقعی خروجی \لر{CLI}، بیشینه تاخیر و افت بسته تکیه دارند.

در \لر{OSPF}، \لر{ping} از \لر{PC1} به \لر{PC3} چهار پاسخ موفق داشت؛
کمینه صفر، بیشینه ۳ و میانگین صفر میلی‌ثانیه ثبت شد. \لر{ping} به
\لر{Internet} نیز چهار پاسخ با زمان کمتر از یک میلی‌ثانیه داشت.

\شروع{شکل}[H]
\centerimg{q04-ospf-pc1-ping-results.png}{12cm}
\شرح{نتیجه \لر{ping} داخلی و \لر{Internet} در \لر{OSPF}}
\برچسب{شکل:q04-ospf-ping}
\پایان{شکل}

در \لر{EIGRP}، هر دو مقصد داخلی و \لر{Internet} بدون افت پاسخ دادند. برای
مقصد داخلی بیشینه یک میلی‌ثانیه و میانگین صفر و برای \لر{Internet} نیز بیشینه
یک و میانگین صفر میلی‌ثانیه ثبت شد.

\شروع{شکل}[H]
\centerimg{q04-eigrp-pc1-ping-results.png}{10cm}
\شرح{نتیجه \لر{ping} داخلی و \لر{Internet} در \لر{EIGRP}}
\برچسب{شکل:q04-eigrp-ping}
\پایان{شکل}

در \لر{RIP}، آزمون کوتاه داخلی بیشینه و میانگین صفر و آزمون \لر{Internet}
بیشینه یک و میانگین صفر میلی‌ثانیه داشت. در دسته بزرگ‌تر، پاسخ نخست از دست رفت
و پاسخ‌های بعدی مقادیر کمتر از یک، ۱۰، یک و ۸ میلی‌ثانیه را نشان دادند؛ این
نوسان اثر صف، آغاز ترجمه و پردازش شبیه‌ساز را بهتر از آزمون چهار بسته‌ای
آشکار کرد.

\شروع{شکل}[H]
\centerimg{q04-rip-pc1-ping-results.png}{10cm}
\شرح{نتیجه \لر{ping} داخلی و \لر{Internet} در \لر{RIP}}
\برچسب{شکل:q04-rip-ping}
\پایان{شکل}

\شروع{شکل}[H]
\centerimg{q04-rip-batch-ping-delay.png}{10cm}
\شرح{بخشی از \لر{ping} دسته‌ای \لر{RIP} و مشاهده نوسان تاخیر}
\برچسب{شکل:q04-rip-batch-ping}
\پایان{شکل}

در مقایسه پایه، هر سه پروتکل ارتباط کامل با میانگین نمایش‌داده‌شده صفر
میلی‌ثانیه داشتند و این دقت زمانی برای رتبه‌بندی قطعی کافی نبود. تفاوت اصلی در
انتخاب مسیر پس از تغییر متریک و در سرعت بازیابی پس از خرابی آشکار شد.

\section{۵. تنظیم \لر{Internet} به‌عنوان مسیر پیش‌فرض}

\لر{R5} مرز شبکه داخلی است و یک مسیر ایستای \لر{0.0.0.0/0} به روتر
\لر{Internet} دارد. \لر{OSPF} و \لر{RIP} آن را با
\کد{default-information originate} و \لر{EIGRP} آن را با
\کد{redistribute static} و \لر{seed metric} به شبکه داخلی معرفی می‌کنند.
در نتیجه مقصد ناشناخته \لر{1.2.3.4} از \لر{PC1} ابتدا به \لر{10.10.6.2}،
سپس \لر{R5} با نشانی \لر{10.10.11.2} و در ادامه روتر \لر{Internet} با
نشانی \لر{213.80.11.5} تحویل داده شد.

\شروع{شکل}[H]
\centerimg{q05-ospf-r5-default-and-internal-interfaces.png}{13cm}
\شرح{مسیر پیش‌فرض \لر{R5} و محدود بودن \لر{OSPF} به دو واسط داخلی}
\برچسب{شکل:q05-ospf-r5-default}
\پایان{شکل}

\شروع{شکل}[H]
\centerimg{q05-ospf-default-route-traceroute.png}{10cm}
\شرح{هدایت مقصد ناشناخته از مسیر پیش‌فرض در فایل \لر{OSPF}}
\برچسب{شکل:q05-ospf-traceroute}
\پایان{شکل}

\شروع{شکل}[H]
\centerimg{q05-eigrp-r5-default-and-internal-interfaces.png}{13cm}
\شرح{مسیر پیش‌فرض \لر{R5} و واسط‌های داخلی \لر{EIGRP}}
\برچسب{شکل:q05-eigrp-r5-default}
\پایان{شکل}

\شروع{شکل}[H]
\centerimg{q05-eigrp-default-route-traceroute.png}{10cm}
\شرح{هدایت مقصد ناشناخته از مسیر پیش‌فرض در فایل \لر{EIGRP}}
\برچسب{شکل:q05-eigrp-traceroute}
\پایان{شکل}

در \لر{RIP}، بخش \لر{running-config} وجود
\کد{default-information originate} را نشان می‌دهد و \لر{traceroute} همان
زنجیره \لر{PC1-R1-R5-Internet} را تایید می‌کند.

\شروع{شکل}[H]
\centerimg{q05-rip-r5-default-origination.png}{10cm}
\شرح{پیکربندی انتشار مسیر پیش‌فرض در \لر{R5} برای \لر{RIP}}
\برچسب{شکل:q05-rip-r5-default}
\پایان{شکل}

\شروع{شکل}[H]
\centerimg{q05-rip-default-route-traceroute.png}{10cm}
\شرح{هدایت مقصد ناشناخته از مسیر پیش‌فرض در فایل \لر{RIP}}
\برچسب{شکل:q05-rip-traceroute}
\پایان{شکل}

\section{۶. جلوگیری از افشای شبکه داخلی برای \لر{Internet}}

هیچ‌یک از سه پروتکل روی واسط بیرونی \لر{R5} اجرا نشد و روی روتر
\لر{Internet} نیز فرایند \لر{OSPF}، \لر{EIGRP} یا \لر{RIP} وجود ندارد.
جدول \لر{Internet} فقط شبکه مستقیم \لر{213.80.11.0/24} و نشانی محلی
\لر{213.80.11.5} را نشان می‌دهد و هیچ مسیر \لر{10.10.0.0/16} در آن دیده
نمی‌شود.

\شروع{شکل}[H]
\centerimg{q06-ospf-internet-routing-isolation.png}{13cm}
\شرح{نبود پروتکل پویا و مسیر داخلی در روتر \لر{Internet}}
\برچسب{شکل:q06-internet-isolation}
\پایان{شکل}

برای بازگشت پاسخ‌ها بدون افزودن مسیر داخلی به \لر{Internet}، روی \لر{R5} از
\لر{NAT overload} استفاده شد. \لر{Fa0/0} و \لر{Fa0/1} نقش \لر{inside} و
\لر{Fa1/0} نقش \لر{outside} دارند. شبکه \لر{10.10.0.0/16} اجازه ترجمه
دارد و نشانی \لر{PC1} یعنی \لر{10.10.6.1} هنگام خروج به
\لر{213.80.11.4} تبدیل می‌شود.

\شروع{شکل}[H]
\centerimg{q06-ospf-r5-nat-and-protocol-boundary.png}{9cm}
\شرح{نقش‌های \لر{NAT} و محدود شدن \لر{OSPF} به واسط‌های داخلی \لر{R5}}
\برچسب{شکل:q06-ospf-nat-boundary}
\پایان{شکل}

\شروع{شکل}[H]
\centerimg{q06-ospf-r5-nat-translations.png}{13cm}
\شرح{ترجمه‌های \لر{ICMP} از \لر{10.10.6.1} به \لر{213.80.11.4} در فایل \لر{OSPF}}
\برچسب{شکل:q06-ospf-nat-translations}
\پایان{شکل}

در فایل \لر{EIGRP}، فهرست همسایه‌های \لر{R5} فقط \لر{R1} و \لر{R4} را
شامل می‌شود و هیچ همسایه‌ای روی \لر{Fa1/0} وجود ندارد. جدول \لر{NAT} نیز
ترجمه فعال و افزایش شمارنده \لر{hit} را نشان می‌دهد.

\شروع{شکل}[H]
\centerimg{q06-eigrp-r5-only-internal-neighbors.png}{12cm}
\شرح{وجود فقط دو همسایه داخلی \لر{EIGRP} در \لر{R5}}
\برچسب{شکل:q06-eigrp-internal-neighbors}
\پایان{شکل}

\شروع{شکل}[H]
\centerimg{q06-eigrp-r5-nat-translations.png}{13cm}
\شرح{ترجمه‌های \لر{ICMP} و آمار \لر{NAT} در فایل \لر{EIGRP}}
\برچسب{شکل:q06-eigrp-nat}
\پایان{شکل}

در فایل \لر{RIP} نیز ترجمه‌های \لر{ICMP} فعال، \لر{inside local} برابر
\لر{10.10.6.1} و \لر{inside global} برابر \لر{213.80.11.4} ثبت شد.
بدین ترتیب پاسخ \لر{Internet} به نشانی عمومی \لر{R5} بازمی‌گردد و نیازی به
یادگیری شبکه داخلی ندارد.

\شروع{شکل}[H]
\centerimg{q06-rip-r5-nat-translations.png}{13cm}
\شرح{ترجمه‌های \لر{ICMP} و آمار \لر{NAT} در فایل \لر{RIP}}
\برچسب{شکل:q06-rip-nat}
\پایان{شکل}

\section{۷. انتقال نیمی از بسته‌ها از \لر{R2} و نیم دیگر از \لر{R5}}

این خواسته دو مرحله مستقل دارد: نخست پروتکل باید دو \لر{next hop} قابل استفاده
برای یک مقصد در \لر{R1} نصب کند و سپس صفحه ارسال
\پانویس{\lr{Forwarding Plane}} باید بسته‌ها را میان آن‌ها تقسیم کند.
\لر{OSPF} و \لر{RIP} فقط مسیرهای با متریک برابر را برای مسیرهای هم‌هزینه
\پانویس{\lr{Equal-Cost Multi-Path}} نصب می‌کنند. در \لر{OSPF} باید مجموع
\لر{cost} دو مسیر و در \لر{RIP} باید تعداد \لر{hop} دو مسیر برابر شود.
\لر{EIGRP} افزون بر مسیرهای هم‌هزینه، با \لر{variance} می‌تواند مسیرهای
نابرابر را نیز در صورت برقراری شرط \لر{Feasible Successor} وارد جدول کند؛
ولی تقسیم دقیق نصف‌به‌نصف به وزن یکسان دو مسیر نیاز دارد.

در \لر{Packet Tracer 8.2.1}، دستور \کد{maximum-paths} روی \لر{IOS}
شبیه‌سازی‌شده \لر{OSPF} پذیرفته نشد و جدول \لر{R1} فقط یک \لر{next hop}
را نگه داشت.

\شروع{شکل}[H]
\centerimg{q07-ospf-packet-tracer-ecmp-limitation.png}{11cm}
\شرح{رد شدن \کد{maximum-paths} و باقی ماندن یک مسیر در \لر{OSPF}}
\برچسب{شکل:q07-ospf-limitation}
\پایان{شکل}

\شروع{شکل}[H]
\centerimg{q07-ospf-r1-route-table.png}{13cm}
\شرح{جدول \لر{R1} و مسیرهای موجود در فایل \لر{OSPF}}
\برچسب{شکل:q07-ospf-r1-routes}
\پایان{شکل}

در \لر{EIGRP}، \لر{variance 1} پذیرفته شد، اما \کد{maximum-paths} و
\کد{traffic-share balanced} در این مدل شبیه‌سازی پشتیبانی نشدند. خروجی
\لر{topology} نیز یک \لر{Successor} را ثبت کرد.

\شروع{شکل}[H]
\centerimg{q07-eigrp-packet-tracer-ecmp-limitation.png}{13cm}
\شرح{محدودیت فرمان‌های توازن بار \لر{EIGRP} در \لر{Packet Tracer}}
\برچسب{شکل:q07-eigrp-limitation}
\پایان{شکل}

\شروع{شکل}[H]
\centerimg{q07-eigrp-r1-route-table.png}{13cm}
\شرح{جدول \لر{R1} و مسیر انتخاب‌شده \لر{EIGRP}}
\برچسب{شکل:q07-eigrp-r1-routes}
\پایان{شکل}

در \لر{RIP} نیز \کد{offset-list} و \کد{maximum-paths} پشتیبانی نشدند و
جدول \لر{R1} یک مسیر فعال را نشان داد.

\شروع{شکل}[H]
\centerimg{q07-rip-packet-tracer-ecmp-limitation.png}{11cm}
\شرح{محدودیت \کد{offset-list} و \کد{maximum-paths} در \لر{RIP}}
\برچسب{شکل:q07-rip-limitation}
\پایان{شکل}

\شروع{شکل}[H]
\centerimg{q07-rip-r1-route-table.png}{13cm}
\شرح{جدول \لر{R1} و مسیر انتخاب‌شده \لر{RIP}}
\برچسب{شکل:q07-rip-r1-routes}
\پایان{شکل}

روی یک روتر فیزیکی \لر{Cisco}، پس از برابر کردن متریک‌ها و نصب دو مسیر،
\لر{CEF} با \کد{ip load-sharing per-packet} روی هر دو واسط خروجی
\لر{Fa0/1} و \لر{Fa1/0} بسته‌های متوالی را به‌صورت \لر{round-robin}
توزیع می‌کند. \کد{maximum-paths 2} تعداد مسیرهای نصب‌شده را به دو محدود
می‌کند. ارسال به‌ازای بسته\پانویس{\lr{Per-Packet Forwarding}} می‌تواند ترتیب
بسته‌ها را تغییر دهد؛ ازاین‌رو در شبکه واقعی معمولاً \لر{ECMP} به‌ازای جریان
یا مقصد ترجیح داده می‌شود و با تعداد زیاد جریان، بار تقریباً متوازن می‌گردد.
نتیجه عملی این بخش آن است که \لر{Packet Tracer} نصب و ارسال دقیق نصف‌به‌نصف
را نمایش نداد و پاسخ کامل آن شامل سازوکار \لر{CEF} روی \لر{IOS} واقعی است.

\section{۸. ترجیح \لر{R3} برای ترافیک غیرمستقیم \لر{R4}}

شبکه‌های مستقیماً متصل \لر{R4} طبیعتاً بدون \لر{next hop} باقی می‌مانند؛
بنابراین خواسته بخش هشتم برای تمام مقصدهای غیرمستقیم تفسیر شد. در \لر{OSPF}،
\لر{cost} واسط \لر{Fa1/0} رو به \لر{R5} روی ۱۰۰۰ قرار گرفت. مسیر \لر{R5}
همچنان به‌عنوان مسیر قابل استفاده باقی ماند، ولی مسیر کم‌هزینه‌تر از \لر{R3}
انتخاب شد.

\شروع{شکل}[H]
\centerimg{q08-ospf-r4-cost-configuration.png}{13cm}
\شرح{افزایش \لر{cost} واسط \لر{R4} به \لر{R5} در \لر{OSPF}}
\برچسب{شکل:q08-ospf-cost}
\پایان{شکل}

\شروع{شکل}[H]
\centerimg{q08-ospf-r4-routes-via-r3.png}{12cm}
\شرح{انتخاب \لر{10.10.10.1} در \لر{R3} برای مسیرهای داخلی \لر{R4}}
\برچسب{شکل:q08-ospf-routes-r3}
\پایان{شکل}

در \لر{EIGRP}، \لر{bandwidth} واسط \لر{R4} به \لر{R5} برابر
\لر{64 Kbit/s} و \لر{delay} برابر ۲۰۰۰۰۰ میکروثانیه شد. این دو مقدار متریک
ترکیبی مسیر \لر{R5} را افزایش دادند و مسیر \لر{R3} با \لر{next hop} برابر
\لر{10.10.10.1} ترجیح داده شد.

\شروع{شکل}[H]
\centerimg{q08-eigrp-r4-bandwidth-delay-configuration.png}{13cm}
\شرح{تغییر \لر{bandwidth} و \لر{delay} واسط \لر{R4} به \لر{R5}}
\برچسب{شکل:q08-eigrp-metric}
\پایان{شکل}

\شروع{شکل}[H]
\centerimg{q08-eigrp-r4-routes-via-r3.png}{12cm}
\شرح{انتخاب مسیرهای داخلی از \لر{R3} در جدول \لر{R4}}
\برچسب{شکل:q08-eigrp-routes-r3}
\پایان{شکل}

در \لر{RIP}، کاهش \لر{bandwidth} یا افزایش \لر{delay} به‌تنهایی انتخاب
مسیر را تغییر نمی‌دهد، زیرا متریک این پروتکل فقط تعداد \لر{hop} است. بنابراین
برخلاف \لر{OSPF} و \لر{EIGRP}، مسیر کم‌سرعت تا زمانی که تعداد \لر{hop} آن
برتر باشد فعال می‌ماند. روش مستقیم افزایش متریک دریافتی از \لر{R5} استفاده
از \کد{offset-list} است؛ این فرمان در \لر{IOS} شبیه‌سازی‌شده پذیرفته نشد.
برای نمایش نتیجه در \لر{Packet Tracer}، مسیرهای مقصد داخلی در \لر{R4} با
\لر{next hop} برابر \لر{10.10.10.1} قرار گرفتند. جدول مسیریابی درستی جهت
خروج به \لر{R3} را نشان می‌دهد، درحالی‌که پیوند \لر{R5} برای مرحله خرابی
بعدی حفظ شد.

\شروع{شکل}[H]
\centerimg{q08-rip-offset-list-limitation.png}{11cm}
\شرح{رد شدن \کد{offset-list} در \لر{RIP} شبیه‌سازی‌شده}
\برچسب{شکل:q08-rip-offset-limitation}
\پایان{شکل}

\شروع{شکل}[H]
\centerimg{q08-rip-r4-static-preference-configuration.png}{9cm}
\شرح{پیکربندی ترجیح مسیرهای داخلی \لر{R4} به سمت \لر{R3}}
\برچسب{شکل:q08-rip-r4-config}
\پایان{شکل}

\شروع{شکل}[H]
\centerimg{q08-rip-r4-routes-via-r3.png}{12cm}
\شرح{نتیجه جدول مسیریابی \لر{R4} پس از اعمال ترجیح \لر{R3}}
\برچسب{شکل:q08-rip-routes-r3}
\پایان{شکل}

\section{۹. واکنش پروتکل‌ها به خاموش شدن \لر{R2}}

برای یکسان بودن نوع خرابی، \لر{R2} در هر سه فایل از کلید \لر{I/O} در زبانه
\لر{Physical} خاموش شد؛ بنابراین هر دو پیوند \لر{R1-R2} و \لر{R2-R3}
هم‌زمان از مدار خارج شدند. پیش از خاموش کردن، از \لر{PC3} فرمان
\کد{ping -n 1000 10.10.6.1} اجرا شد و پس از بازیابی مسیر با \لر{tracert}
جهت جدید کنترل گردید.

\شروع{شکل}[H]
\centerimg{q09-r2-physical-power-switch.png}{11cm}
\شرح{کلید فیزیکی مورد استفاده برای خاموش کردن کامل \لر{R2}}
\برچسب{شکل:q09-r2-power}
\پایان{شکل}

\subsection{واکنش \لر{OSPF}}

در \لر{OSPF}، پیش از خرابی مسیر \لر{PC1} از \لر{R3} و مسیر پیش‌فرض از
\لر{R5} ثبت شده بود. هنگام تغییر مسیر، از ۲۴ بسته یک بسته از دست رفت و سپس
\لر{TTL} از ۱۲۴ به ۱۲۵ تغییر کرد که کوتاه‌تر شدن مسیر را نشان می‌دهد.
\لر{tracert} نهایی مسیر \لر{PC3-R4-R5-R1-PC1} را با \لر{hop}های
\لر{10.10.16.2}، \لر{10.10.12.2}، \لر{10.10.11.1} و \لر{10.10.6.1}
تایید کرد.

\شروع{شکل}[H]
\centerimg{q09-ospf-r4-routes-before-r2-failure.png}{12cm}
\شرح{مسیرهای \لر{R4} پیش از خاموش شدن \لر{R2} در فایل \لر{OSPF}}
\برچسب{شکل:q09-ospf-before}
\پایان{شکل}

\شروع{شکل}[H]
\centerimg{q09-ospf-ping-during-r2-failure.png}{10cm}
\شرح{یک \لر{timeout} و بازیابی \لر{ping} در \لر{OSPF}}
\برچسب{شکل:q09-ospf-ping}
\پایان{شکل}

\شروع{شکل}[H]
\centerimg{q09-ospf-traceroute-after-r2-failure.png}{10cm}
\شرح{مسیر بازیابی‌شده \لر{OSPF} از \لر{R5}}
\برچسب{شکل:q09-ospf-trace}
\پایان{شکل}

\subsection{واکنش \لر{EIGRP}}

در \لر{EIGRP}، \لر{topology} پیش از خرابی یک \لر{Successor} از \لر{R3}
را نشان می‌دهد. پس از قطع \لر{R2}، \لر{DUAL} مسیر را دوباره محاسبه کرد و هر
۱۵ بسته ثبت‌شده بدون افت پاسخ گرفتند. تغییر \لر{TTL} از ۱۲۴ به ۱۲۵ و
\لر{tracert} نهایی عبور از \لر{R5} را تایید کردند.

\شروع{شکل}[H]
\centerimg{q09-eigrp-r4-routes-before-r2-failure.png}{11cm}
\شرح{مسیرهای \لر{R4} پیش از خرابی در فایل \لر{EIGRP}}
\برچسب{شکل:q09-eigrp-before}
\پایان{شکل}

\شروع{شکل}[H]
\centerimg{q09-eigrp-topology-before-r2-failure.png}{12cm}
\شرح{\لر{Successor} و متریک مسیر \لر{PC1} پیش از خرابی}
\برچسب{شکل:q09-eigrp-topology}
\پایان{شکل}

\شروع{شکل}[H]
\centerimg{q09-eigrp-ping-during-r2-failure.png}{10cm}
\شرح{بازیابی \لر{EIGRP} بدون افت در نمونه ثبت‌شده}
\برچسب{شکل:q09-eigrp-ping}
\پایان{شکل}

\شروع{شکل}[H]
\centerimg{q09-eigrp-traceroute-after-r2-failure.png}{10cm}
\شرح{مسیر بازیابی‌شده \لر{EIGRP} از \لر{R5}}
\برچسب{شکل:q09-eigrp-trace}
\پایان{شکل}

\subsection{واکنش \لر{RIP}}

در \لر{RIP}، مسیر فعال پیش از خرابی به سمت \لر{R3} ثبت شده بود. در نمونه
ثبت‌شده، هر ۲۲ بسته پاسخ گرفت و تغییر \لر{TTL} گذار به مسیر کوتاه‌تر
\لر{R5} را نشان داد. \لر{tracert} نیز همان مسیر \لر{R4-R5-R1} را ثبت کرد.
\لر{RIP} در حالت عمومی برای انتشار دوره‌ای و زمان‌سنج‌های ۳۰، ۱۸۰ و ۲۴۰
ثانیه‌ای به زمان بیشتری نیاز دارد؛ خاموش شدن فیزیکی پیوند و آماده بودن مسیر
جایگزین در این اجرای کوچک باعث شد وقفه در نمونه \لر{ping} دیده نشود.

\شروع{شکل}[H]
\centerimg{q09-rip-r4-routes-before-r2-failure.png}{11cm}
\شرح{مسیرهای \لر{R4} پیش از خرابی در فایل \لر{RIP}}
\برچسب{شکل:q09-rip-before}
\پایان{شکل}

\شروع{شکل}[H]
\centerimg{q09-rip-ping-during-r2-failure.png}{10cm}
\شرح{گذار مسیر \لر{RIP} بدون افت در نمونه ثبت‌شده}
\برچسب{شکل:q09-rip-ping}
\پایان{شکل}

\شروع{شکل}[H]
\centerimg{q09-rip-traceroute-after-r2-failure.png}{10cm}
\شرح{مسیر بازیابی‌شده \لر{RIP} از \لر{R5}}
\برچسب{شکل:q09-rip-trace}
\پایان{شکل}

از نظر سازوکار، \لر{EIGRP} با \لر{DUAL} و استفاده فوری از
\لر{Feasible Successor} در صورت وجود، سریع‌ترین واکنش را ارائه می‌کند.
\لر{OSPF} پس از انتشار \لر{LSA} و اجرای دوباره \لر{SPF} همگرایی سریعی دارد،
اما در این ثبت یک بسته از دست رفت. \لر{RIP} ذاتاً به به‌روزرسانی بردار فاصله
\پانویس{\lr{Distance Vector}} و زمان‌سنج‌ها وابسته است و در شبکه‌های بزرگ‌تر
کندترین همگرایی را دارد. در داده‌های این اجرا، \لر{EIGRP} و \لر{RIP} بدون
افت ثبت شدند و \لر{OSPF} یک افت داشت؛ بااین‌حال نتیجه نظری و مقیاس‌پذیر
همچنان ترتیب \لر{EIGRP}، سپس \لر{OSPF} و در نهایت \لر{RIP} است.
